\documentclass[10pt,a4paper]{amsart}
\usepackage{amsmath,amssymb,amsthm,graphicx}
\usepackage{verbatim}
\newenvironment{reference}{\comment}{\endcomment}
%\newenvironment{reference}{}{}
\newenvironment{slogan}{\comment}{\endcomment}
\newenvironment{history}{\comment}{\endcomment}

% For commutative diagrams we use Xy-pic
\usepackage[all]{xy}

% We use 2cell for 2-commutative diagrams.
\xyoption{2cell}
% Theorem environments.
%
\theoremstyle{plain}
\newtheorem{theorem}[subsection]{Theorem}
\newtheorem{proposition}[subsection]{Proposition}
\newtheorem{lemma}[subsection]{Lemma}

\theoremstyle{definition}
\newtheorem{definition}[subsection]{Definition}
\newtheorem{example}[subsection]{Example}
\newtheorem{exercise}[subsection]{Exercise}
\newtheorem{situation}[subsection]{Situation}

\theoremstyle{remark}
\newtheorem{remark}[subsection]{Remark}
\newtheorem{remarks}[subsection]{Remarks}

\numberwithin{equation}{subsection}

% Macros
%
\def\lim{\mathop{\mathrm{lim}}\nolimits}
\def\colim{\mathop{\mathrm{colim}}\nolimits}
\def\Spec{\mathop{\mathrm{Spec}}}
\def\Hom{\mathop{\mathrm{Hom}}\nolimits}
\def\Ext{\mathop{\mathrm{Ext}}\nolimits}
\def\SheafHom{\mathop{\mathcal{H}\!\mathit{om}}\nolimits}
\def\SheafExt{\mathop{\mathcal{E}\!\mathit{xt}}\nolimits}
\def\Sch{\mathit{Sch}}
\def\Mor{\mathop{\mathrm{Mor}}\nolimits}
\def\Ob{\mathop{\mathrm{Ob}}\nolimits}
\def\Sh{\mathop{\mathit{Sh}}\nolimits}
\def\NL{\mathop{N\!L}\nolimits}
\def\CH{\mathop{\mathrm{CH}}\nolimits}
\def\proetale{{pro\text{-}\acute{e}tale}}
\def\etale{{\acute{e}tale}}
\def\QCoh{\mathit{QCoh}}
\def\Ker{\mathop{\mathrm{Ker}}}
\def\Im{\mathop{\mathrm{Im}}}
\def\Coker{\mathop{\mathrm{Coker}}}
\def\Coim{\mathop{\mathrm{Coim}}}

% Boxtimes
%
\DeclareMathSymbol{\boxtimes}{\mathbin}{AMSa}{"02}

%
% Macros for moduli stacks/spaces
%
\def\QCohstack{\mathcal{QC}\!\mathit{oh}}
\def\Cohstack{\mathcal{C}\!\mathit{oh}}
\def\Spacesstack{\mathcal{S}\!\mathit{paces}}
\def\Quotfunctor{\mathrm{Quot}}
\def\Hilbfunctor{\mathrm{Hilb}}
\def\Curvesstack{\mathcal{C}\!\mathit{urves}}
\def\Polarizedstack{\mathcal{P}\!\mathit{olarized}}
\def\Complexesstack{\mathcal{C}\!\mathit{omplexes}}
% \Pic is the operator that assigns to X its picard group, usage \Pic(X)
% \Picardstack_{X/B} denotes the Picard stack of X over B
% \Picardfunctor_{X/B} denotes the Picard functor of X over B
\def\Pic{\mathop{\mathrm{Pic}}\nolimits}
\def\Picardstack{\mathcal{P}\!\mathit{ic}}
\def\Picardfunctor{\mathrm{Pic}}
\def\Deformationcategory{\mathcal{D}\!\mathit{ef}}
\usepackage{cleveref}
\crefname{originalreference}{original reference}{original references}
\Crefname{originalreference}{Original reference}{Original references}
\makeatletter
\newlabel{lemma-Gm-torsor}{{0EQG}{source}{}{sourcepointer-1}{}}
\newlabel{lemma-Gm-torsor@cref}{{[originalreference][0][]0EQG}{[1][0][]source}{}{}{}}
\makeatother
\crefname{originalreference}{original reference}{original references}
\Crefname{originalreference}{Original reference}{Original references}
\makeatletter
\newlabel{lemma-Gm-torsor-divisor-meromorphic-section}{{0EQH}{source}{}{sourcepointer-2}{}}
\newlabel{lemma-Gm-torsor-divisor-meromorphic-section@cref}{{[originalreference][0][]0EQH}{[1][0][]source}{}{}{}}
\makeatother
\crefname{originalreference}{original reference}{original references}
\Crefname{originalreference}{Original reference}{Original references}
\makeatletter
\newlabel{lemma-compose-flat-pullback}{{02RD}{source}{}{sourcepointer-3}{}}
\newlabel{lemma-compose-flat-pullback@cref}{{[originalreference][0][]02RD}{[1][0][]source}{}{}{}}
\makeatother
\crefname{originalreference}{original reference}{original references}
\Crefname{originalreference}{Original reference}{Original references}
\makeatletter
\newlabel{lemma-prepare-flat-pullback-rational-equivalence}{{0EPH}{source}{}{sourcepointer-4}{}}
\newlabel{lemma-prepare-flat-pullback-rational-equivalence@cref}{{[originalreference][0][]0EPH}{[1][0][]source}{}{}{}}
\makeatother
\crefname{originalreference}{original reference}{original references}
\Crefname{originalreference}{Original reference}{Original references}
\makeatletter
\newlabel{situation-setup}{{02QL}{source}{}{sourcepointer-5}{}}
\newlabel{situation-setup@cref}{{[originalreference][0][]02QL}{[1][0][]source}{}{}{}}
\makeatother
\title{Flat pullback of the divisor of a meromorphic section is the multiplicity weighted sum over the components of the source}
\author{The Stacks Project Authors}
\date{stacks-project:spaces-chow.tex:local-snapshot-20261012; GNU FDL 1.2 or later}
\begin{document}
\maketitle
\noindent\textit{Editorial note.} All delivered mathematical text is the author's own, copied from the pinned Stacks Project chapter file spaces-chow.tex; no statement or proof step was rewritten, reordered or removed.\par\medskip
\noindent\textit{Editorial note.} The counted proof is the author's complete argument: a reduction to the trivialising torsor followed by component and multiplicity bookkeeping and the identification of the sections with rational functions on the torsor\par\medskip
\noindent\textit{Editorial note.} The label lemma-prepare-flat-pullback-cap-c1 names the counted statement, and the proof environment immediately following it in the source is delivered verbatim as its proof.\par\medskip
\noindent\textit{Editorial note.} The counted statement is carried by the author's own lemma environment in the chapter source.\par\medskip
\noindent\textit{Allowed standard prerequisites (editorial).}\begin{itemize}
\item The divisor of a nonzero meromorphic section of an invertible sheaf and its flat pullback, together with the multiplicity of an irreducible component in a cycle
\item The construction of a torsor trivialising an invertible sheaf, the identification of the divisor of a meromorphic section with the divisor of the corresponding rational function on that torsor, and the behaviour of the divisors under the projection
\item The composition of flat pullbacks, and the rational function case of the formula for flat pullback of a principal divisor with multiplicities
\end{itemize}\medskip
\ifdefined\hypertarget\hypertarget{sourcepointer-1}{}\fi
\noindent\textit{Source reference.} Original reference 0EQG: sources/stacks-spaces-chow/pinned-original-source.tex, source lines 2385--2385. Prerequisite pointer invoked at proof line(s) 2757; it names a standard fact of the same chapter that the counted argument uses, and every step of the counted proof built on it is retained verbatim in the delivered body.\par\medskip
\ifdefined\hypertarget\hypertarget{sourcepointer-2}{}\fi
\noindent\textit{Source reference.} Original reference 0EQH: sources/stacks-spaces-chow/pinned-original-source.tex, source lines 2477--2477. Prerequisite pointer invoked at proof line(s) 2778; it names a standard fact of the same chapter that the counted argument uses, and every step of the counted proof built on it is retained verbatim in the delivered body.\par\medskip
\ifdefined\hypertarget\hypertarget{sourcepointer-3}{}\fi
\noindent\textit{Source reference.} Original reference 02RD: sources/stacks-spaces-chow/pinned-original-source.tex, source lines 1347--1347. Prerequisite pointer invoked at proof line(s) 2779; it names a standard fact of the same chapter that the counted argument uses, and every step of the counted proof built on it is retained verbatim in the delivered body.\par\medskip
\ifdefined\hypertarget\hypertarget{sourcepointer-4}{}\fi
\noindent\textit{Source reference.} Original reference 0EPH: sources/stacks-spaces-chow/pinned-original-source.tex, source lines 2039--2039. Prerequisite pointer invoked at proof line(s) 2753; it names a standard fact of the same chapter that the counted argument uses, and every step of the counted proof built on it is retained verbatim in the delivered body.\par\medskip
\ifdefined\hypertarget\hypertarget{sourcepointer-5}{}\fi
\noindent\textit{Source reference.} Original reference 02QL: sources/stacks-spaces-chow/pinned-original-source.tex, source lines 264--264. Prerequisite pointer invoked at proof line(s) 2736; it names a standard fact of the same chapter that the counted argument uses, and every step of the counted proof built on it is retained verbatim in the delivered body.\par\medskip
\begin{lemma}
\label{lemma-prepare-flat-pullback-cap-c1}
In Situation \ref{situation-setup} let $X, Y/B$ be good.
Let $f : X \to Y$ be a flat morphism of relative dimension $r$.
Let $\mathcal{L}$ be an invertible sheaf on $Y$.
Assume $Y$ is integral and $n = \dim_\delta(Y)$.
Let $s$ be a nonzero meromorphic section of $\mathcal{L}$.
Then we have
$$
f^*\text{div}_\mathcal{L}(s) = \sum n_i\text{div}_{f^*\mathcal{L}|_{X_i}}(s_i)
$$
in $Z_{n + r - 1}(X)$. Here the sum is over the irreducible
components $X_i \subset X$ of $\delta$-dimension $n + r$,
the section $s_i = f|_{X_i}^*(s)$ is the pullback of $s$, and
$n_i = m_{X_i, X}$ is the multiplicity of $X_i$ in $X$.
\end{lemma}

\begin{proof}
Using sleight of hand we will deduce this from
Lemma \ref{lemma-prepare-flat-pullback-rational-equivalence}.
(An alternative is to redo the proof of that lemma in the
setting of meromorphic sections of invertible modules.)
Namely, let $q : T \to Y$ be the morphism of
Lemma \ref{lemma-Gm-torsor} constructed using $\mathcal{L}$ on $Y$.
We will use all the properties of $T$ stated in this lemma.
Consider the fibre product diagram
$$
\xymatrix{
T' \ar[r]_{q'} \ar[d]_h & X \ar[d]^f \\
T \ar[r]^q & Y
}
$$
Then $q' : T' \to X$ is the morphism constructed using $f^*\mathcal{L}$
on $X$. Then it suffices to prove
$$
(q')^*f^*\text{div}_\mathcal{L}(s) =
\sum n_i (q')^*\text{div}_{f^*\mathcal{L}|_{X_i}}(s_i)
$$
Observe that $T'_i = q^{-1}(X_i)$ are the irreducible components of $T'$
and that $n_i$ is the multiplicity of $T'_i$ in $T'$.
The left hand side is equal to
$$
h^*q^*\text{div}_\mathcal{L}(s) = h^*\text{div}_T(q^*(s))
$$
by Lemma \ref{lemma-Gm-torsor-divisor-meromorphic-section}
(and Lemma \ref{lemma-compose-flat-pullback}).
On the other hand, denoting $q'_i : T'_i \to X_i$
the restriction of $q'$ we find that
Lemma \ref{lemma-Gm-torsor-divisor-meromorphic-section}
also tells us the right hand side is equal to
$$
\sum n_i \text{div}_{T_i}((q'_i)^*(s_i))
$$
In these two formulas the expressions $q^*(s)$ and $(q'_i)^*(s_i)$
represent the rational functions corresponding to the pulled back
meromorphic sections of $q^*\mathcal{L}$ and $(q'_i)^*f^*\mathcal{L}|_{X_i}$
via the isomorphism $\alpha : q^*\mathcal{L} \to \mathcal{O}_T$
and its pullbacks to spaces over $T$. With this convention it is
clear that $(q'_i)^*(s_i)$ is the composition of the rational function
$q^*(s)$ on $T$ and the morphism $h|_{T'_i} : T'_i \to T$.
Thus Lemma \ref{lemma-prepare-flat-pullback-rational-equivalence}
exactly says that
$$
h^*\text{div}_T(q^*(s)) = \sum n_i \text{div}_{T_i}((q'_i)^*(s_i))
$$
as desired.
\end{proof}

\end{document}
